\documentclass[11pt]{article}

\usepackage[letterpaper,margin=0.5in]{geometry}
\usepackage[scaled]{helvet}
\renewcommand{\familydefault}{\sfdefault}
\usepackage[T1]{fontenc}
\usepackage{microtype}
\usepackage{amsmath,amssymb}
\usepackage{parskip}
\usepackage{xcolor}
\usepackage[hidelinks]{hyperref}
\usepackage{ulem}
\normalem

\setlength{\parskip}{4pt}
\setlength{\parindent}{0pt}
\pagestyle{empty}

\begin{document}

\textbf{Background.} Cellular decisions in development, immunity and
disease are made under stringent physical constraints. In a foundational
analysis of the \textit{Drosophila} Bicoid morphogen gradient, Gregor
\textit{et al.} (\textit{Cell}, 2007) showed that the early embryo
positions the \textit{hunchback} expression boundary at $\sim$1\%
egg-length precision, saturating the Berg--Purcell diffusion-limited bound
on concentration sensing within the available 3-min nuclear-cycle-14
window. The readout machinery sits at the physical floor. This finding is
deep, but it leaves out everything between transcription factor (TF)
binding and a productive transcriptional burst. The ``receptor'' for
transcriptional output is not a fixed binding site but a productive
enhancer--promoter (EP) contact, and EP contact statistics are set by
chromatin polymer dynamics, not by TF diffusion. Recent quantitative
measurements have made the gap concrete. \textbf{Chen et al.} (Gregor
lab, \textit{Nat. Genet.}, 2018) tracked EP topology and transcription
simultaneously at the \textit{eve} locus and identified three distinct
states with very different timescales: open/distal (the bulk),
inactive-paired ($\sim$7-min dwell), and active-paired ($\sim$90-min
dwell). The transition into the paired state is roughly eight-fold slower
than naive polymer first-passage predicts; transcription itself stabilizes
the paired state by an order of magnitude; and a competing reporter
nearby produces adult patterning defects. \textbf{Gabriele et al.}
(Hansen lab, \textit{Science}, 2022) measured CTCF/cohesin loop lifetimes
of tens of minutes in mESCs, occupied only a few percent of the cell
cycle. \textbf{Saiz et al.} (de Wit lab, \textit{bioRxiv}, 2026) used
auxin-degron CTCF depletion in mouse gastruloids to show that fate
specification is preserved without CTCF looping but morphogenesis fails,
and post-fate fine-tuning of expression collapses. \textbf{Lammers et
al.} (Garcia lab, \textit{PNAS}, 2023) derived an information-rate bound
on transcriptional decisions limited by a precision/sharpness/specificity
tradeoff against noncognate-factor interference, with fly and mouse on
opposite sides of this kinetic axis. Together, these results imply that
mammalian fate decisions are constrained by three physically independent
sources of noise --- (i) Berg--Purcell counting of TF molecules,
(ii) the Lammers kinetic-proofreading tradeoff, and (iii) a chromatin
polymer first-passage and dwell-time tradeoff that determines whether
productive EP contacts occur often enough \emph{and} last long enough for
transcription to fire. \textbf{Which of these is rate-limiting in
mammalian development is not known.}

\textbf{Innovation.} \uline{The overall vision of the Upadhyaya lab is to
combine single-molecule imaging, quantitative image analysis, and
statistical-physics-based modeling with cell-biological and genomic
approaches to elucidate the biophysical principles underlying cellular
decision-making.} During the current MIRA cycle we extended this approach
from immune mechanobiology to TF/chromatin dynamics, developing
single-molecule tracking (SMT) and analysis methods for the glucocorticoid
receptor (GR), estrogen receptor (ER), the pioneer factor FOXA1, and the
chromatin remodeler BRG1, and showed that long-residence TF events on
chromatin are quantitatively associated with productive transcription
(Mehta \textit{et al.}, 2022; Garcia \textit{et al.}, 2021). We further
showed that ERs in MCF-7 cells exhibit substrate-stiffness-dependent
shifts in chromatin-bound dynamics
(H2B $\alpha\!\approx\!0.41$ on soft, ER $\alpha\!\approx\!0.61$),
demonstrating that mechanical cues modulate TF/chromatin engagement at
the level of single binding events. \textit{The natural next question is
whether these single-molecule observables --- chromatin
$\alpha$, free fractions, residence-time distributions --- are sufficient
to populate the noise budget that constrains a cell's fate decision in a
true developmental context.} Our laboratory is uniquely positioned to
answer this question: we have an established \textbf{custom HILO
microscope build}, a \textbf{validated SMT analysis stack}
(\textit{smt-pipeline}: fractional Brownian fitting, pEMv2 with bootstrap
CIs, state-array posterior decomposition, two-condition statistical
comparison) released as open source, and a \textbf{new widefield HILO
microscope} dedicated to live SMT in 3D-cultured tissue, currently in
final stages of construction.

\textbf{The 2D mouse gastruloid as the entry point.} The gastruloid grown
on micropatterned (CYTOO) substrates is the cleanest mammalian system in
which the polymer-revised Berg--Purcell bound can be tested by SMT. It
begins as a uniform aggregate of mouse embryonic stem cells, breaks
symmetry under defined Wnt stimulation, and acquires a Brachyury-positive
pole opposite a Sox2-retaining domain over a tractable 24--72-hour
window. It is geometrically defined, optically accessible, reproducible
across replicates, and --- unlike the \textit{in vivo} mouse embryo ---
amenable to single-molecule imaging in living tissue. Every term in the
noise budget except the EP first-passage time itself is observable
through SMT of histone H2B (chromatin) and of fate-relevant TFs (Sox2,
Brachyury). H2B SMT yields the chromatin MSD exponent $\alpha$ and bounds
the polymer first-passage timescale; Sox2 and Bra SMT yield free
fractions, residence-time distributions, and the TF-side prefactors of
the Berg--Purcell and Lammers terms; spatial mapping of these observables
across the polarized axis decomposes the readout into source,
propagation, and readout components and discriminates reaction--diffusion
from localized-source mechanisms of axial symmetry breaking. We are
explicitly \emph{not} proposing live two-color enhancer/promoter pair
imaging, which requires engineered loci and orthogonal labeling chemistry
more naturally pursued by the Gregor, Garcia, or Hansen laboratories;
\textit{our contribution is the parallel single-molecule measurement of
the chromatin and TF dynamics that flank the EP contact event ---
populating noise-budget terms that two-color locus tracking does not
resolve, and providing the quantitative calibration against which
locus-resolved experiments in collaborating laboratories can be
interpreted.}

\textbf{Impact.} A direct single-molecule decomposition of the physical
noise budget that constrains mammalian fate decisions has not been done.
The result will reveal which physical bottleneck is rate-limiting in
mammalian gastrulation, distinguish reaction--diffusion from
localized-source axial symmetry breaking, and clarify the role of
chromatin organization in fate specification independent of its later
role in long-term gene-expression homeostasis. The framework and methods
established here generalize beyond development: the same noise budget
constrains hormone-driven transcriptional decisions in cancer (extending
our ongoing ER work) and antigen-driven fate decisions in lymphocytes
(extending our T-cell program), unifying the lab's three programs under
a common biophysical principle.

\end{document}
